\documentclass[10pt,a4paper]{amsart}
\usepackage[margin=19mm]{geometry}
\usepackage{amsmath,amssymb,mathtools}
\usepackage[scr=rsfs,cal=euler]{mathalfa}
\usepackage{xfrac}
\usepackage[unicode,hidelinks,bookmarks=false]{hyperref}
\newcommand{\ie}{i.e.\ }
\newcommand{\cf}{cf.\ }
\newcommand{\resp}{respectively\ }
\newcommand{\Subsection}{\subsection}
\providecommand{\nobreakdash}{\nobreak\hbox{-}}
\setlength{\emergencystretch}{2em}
\multlinegap0pt
\usepackage[shortlabels]{enumitem}
\usepackage{mathtools}
\usepackage{algorithm}\usepackage{algpseudocode}
\usepackage{bm}
\usepackage{amssymb}
\usepackage{xcolor}
\usepackage{cancel}
\usepackage{tikz}
\usepackage[normalem]{ulem}
\newtheorem{assumption}{Assumption}
\newtheorem{defn}{Definition}
\newtheorem{lemma}{Lemma}
\newcommand{\EE}{\mathbb{E}}
\newcommand{\PP}{\mathbb{P}}
\newcommand{\RR}{\mathbb{R}}
\title{Multistage Conditional Compositional Optimization}
\author{Buse \c{S}en, Yifan Hu, Daniel Kuhn}
\date{Source: arXiv:2604.14075v1 (CC BY 4.0)}
\begin{document}
\maketitle

\noindent\emph{Source and licence.} Multistage Conditional Compositional Optimization. Version arXiv:2604.14075v1 (CC BY 4.0), distributed by arXiv under the Creative Commons Attribution 4.0 International licence, http://creativecommons.org/licenses/by/4.0/, as recorded in the arXiv licence metadata for this exact version and shown on the abstract page https://arxiv.org/abs/2604.14075v1. Full licence notice: \texttt{licenses/arxiv-2604.14075-CC-BY-4.0.txt}.\par\smallskip

\noindent\textit{Editorial scope.} Counted unit: the article's lemma lemma:bias\_bound\_rtmlmc\_grad (source0.tex lines 976--981). The article's complete original proof of it is delivered with it (source0.tex lines 2116--2201, one \texttt{proof} environment of 86 lines, which the article places away from the statement). No mathematics is written, repaired or re-derived: the statement and the proof are the article's own lines, in the article's own notation, and no retained line is substituted or re-flowed.  Retained uncounted as the source-local context the proof needs: lines 382--388; lines 392--396; lines 398--402; lines 673--676; lines 682--684; lines 891--893; lines 905--923; lines 942--944; lines 950--952; lines 1657--1666.  Nothing was omitted from any retained span, so the package carries no omission record.  No retained line contains a citation, so the wrapper carries no bibliography and no bibliography entry.  No retained line contains a figure environment, an image-inclusion command or drawing code, so no image is delivered and the package declares no asset dependency.  Editorial formatting only, in addition: an AMS \texttt{amsart} preamble that restates the article's own declarations and macros used by the retained lines, namely the environments \texttt{assumption}, \texttt{defn}, \texttt{lemma} on separate counters, together with the standard packages those lines need.  The retained environments print in this document as Assumption 1, Definition 1, Lemma 1 in order of appearance, whereas the article prints the same items with the numbering of its own full document.  The complete native source of the article as distributed in the arXiv e-print for this exact version is bundled unchanged as \texttt{sources/198516/source0.tex}.
\begin{assumption}[Basic Regularity Conditions] 
\label{assump:general}
\leavevmode\begin{enumerate}[label=(\roman*)]
    \item The integrand $f_t(\xi_t,x_t)$ is Borel-measurable in~$\xi_t$, and $\EE[\|f_t(\xi_t,0)\|_2]<\infty$ for every $t\in[T]$. 
    \item The feasible set $\mathcal{X}$ is convex and compact. Thus, it has a finite diameter~$D_\mathcal{X}$.
\end{enumerate}
\end{assumption}

\begin{assumption}[Lipschitz Continuity] 
\label{assump:lipschitz}
   $f_t(\xi_t,x_t)$ is $L_t$-Lipschitz continuous in $x_t$ uniformly for all $\xi_t$, i.e.,
    $$ \|f_t(\xi_t,x_t) -  f_t(\xi_t,y_t)\|_2 \leq L_t\|x_t-y_t\|_2 \quad \forall x_t,y_t \in \RR^{d_t},\; \forall \xi_t\in\RR^{m_t}.$$
\end{assumption}

\begin{assumption}[Smoothness] 
\label{assump:smooth}
 $f_t(\xi_t,x_t)$ is $S_t$-smooth in $x_t$ for across all $\xi_{t}$, i.e., 
  $$\|\nabla_{x_t} f_t(\xi_t,x_t) - \nabla_{x_t} f_t(\xi_t,y_t) \|_2\leq S_t\|x_t-y_t\|_2\quad \forall x_t,y_t\in \RR^{d_t},\; \forall \xi_t\in\RR^{m_t}.$$
\end{assumption}

\begin{align}
\label{eq:mu}
    \mu_t^{p}(x) \coloneqq \EE_{t-1}\big[\|\hat{H}_t(x)\|_p^p \big],
\end{align}

\begin{assumption}\label{assump:mu_bar}
    The conditional moment bound $\overline{\mu}_T^p $ of the integrand $f_T(\xi_T,x)$ is finite up to order $p=2$ if Assumption~\ref{assump:smooth} does not hold and up to order $p=2^T$ if Assumption~\ref{assump:smooth} holds.
\end{assumption}

\begin{align}\label{eq:G_t}
    G_t(x) \coloneqq \EE_{t-1}\left[G_{t+1}(x) \nabla_{x_t} f_t(\xi_t, F_{t+1}^T(x))\right]\in\RR^{d\times d_{t-1}}\quad \forall t\in[T-1]
\end{align}

\begin{defn}[MLMC Gradient Estimator]
\label{def:rumlmc-grad}
For any $t\in[T-1]$, select a rate parameter~$r_t \in (0,1)$ and truncation point $M_t\in\mathbb N$, and let $\lambda_t$ be a random variable that follows the truncated geometric distribution~$\mathrm{Geo}(r_t | M_t)$. Assume that~$\lambda_t$ is independent of all other random objects, and set $q_t(\ell)\coloneqq\PP(\lambda_t=\ell)$ for~$\ell=0,\ldots,M_t$ and~$t\in[T-1]$. We construct gradient estimators $\hat G_t(x)$, $t\in[T]$, using the following backward recursion. We initialize the recursion by $\hat G_T(x)\coloneqq\nabla_x f_T(\xi_T,x)$. Next, for any $t\in[T-1]$, we~set
\begin{align}
    \hat{G}_t(x) \coloneqq \frac{1}{q_{t}(\lambda_{t})}\left( \hat{g}_t^{\lambda_t}(x)-\frac{1}{2}\hat{g}_t^{\lambda_t, {\rm e}}(x)-\frac{1}{2}\hat{g}_t^{\lambda_t, {\rm o}}(x)\label{eq:G_hat}\right),
\end{align}
where
\begin{equation*}
    \hat{g}_t^{\ell}(x) \coloneqq \hat\EE_t^{\ell}[\hat{G}_{t+1}(x) ]\; \nabla_{x_t} f_t \big(\xi_t, \hat\EE_t^{\ell}[ \hat{H}_{t+1}(x)] \big)\quad \forall \ell=0,\ldots,M_t
\end{equation*}
as well as
\begin{equation*}
\begin{aligned}
     \hat{g}_t^{\ell, {\rm e}}(x) \coloneqq& \hat\EE_t^{\ell, {\rm e}}[ \hat{G}_{t+1}(x)] \; \nabla_{x_t} f_t \big(\xi_t, \hat\EE_t^{\ell, {\rm e}}[\hat{H}_{t+1}(x)] \big)\quad \forall \ell\in[M_t], \\
     \hat{g}_t^{\ell, {\rm o}}(x) \coloneqq& \hat\EE_t^{\ell, {\rm o}}[ \hat{G}_{t+1}(x)] \; \nabla_{x_t} f_t \big(\xi_t, \hat\EE_t^{\ell, {\rm o}}[\hat{H}_{t+1}(x)] \big)\quad \forall \ell\in[M_t].
\end{aligned}
\end{equation*}
We also set $\hat{g}_t^{0,{\rm e}}(x)\coloneqq\hat{g}_t^{0,{\rm o}}(x)\coloneqq0$. The MLMC estimator for $\nabla F(x)$ is then defined as $ \hat G(x)\coloneqq\hat\EE [\hat{G}_{1}(x)]$. 
\end{defn}

\begin{align}\label{eq:nu}
    \nu_t^{p}(x) \coloneqq \EE_{t-1}\!\left[ \|\hat{G}_t(x)\|_2^p \right].
\end{align}

\begin{assumption}\label{assump:nu_bar}
    The conditional moment bound $\overline{\mu}_T^p $ of the integrand $ f_T(\xi_T,x)$ is finite up to order $p=2^T(\rho_{T-1}+1)$, and the conditional moment bound $\overline{\nu}_T^p $ of the gradient $\nabla f_T(\xi_T,x)$ is finite up to order $p=2^T$.
\end{assumption}

\begin{lemma}
\label{lemma:bias_bound_rtmlmc_H_t} 
If Assumptions~\ref{assump:general}, \ref{assump:lipschitz} and~\ref{assump:mu_bar} hold and if $t\in[T]$, then we have
\begin{align*}
    \big\|\EE_{t-1}[\hat{H}_t(x)] - F_t^T(x)\big\|_2 \leq \begin{cases}
        \sum_{s=t}^{T-1}\frac{L_{[t:s]}\;\EE_{t-1}[\mu_{s+1}^2(x)]^\frac{1}{2}}{2^{M_s/2}}&\text{if Assumption~\ref{assump:smooth} does not hold,}\\
         \sum_{s=t}^{T-1}\frac{L_{[t:s-1]}S_s\;\EE_{t-1}[\mu_{s+1}^{2}(x)]}{2^{M_s + 1}} &\text{if Assumption~\ref{assump:smooth} holds.}
    \end{cases}
\end{align*}
\end{lemma}

\begin{lemma}\label{lemma:bias_bound_rtmlmc_grad}
    If Assumptions~\ref{assump:general}, \ref{assump:lipschitz}, \ref{assump:smooth} and~\ref{assump:nu_bar} hold, the bias of the MLMC gradient estimator satisfies
    \begin{align*}
        \big\| \EE[\hat{G}(x)]- G_1(x) \big\|_2 \leq \sum_{t=1}^{T-1}L_{[t-1]}S_t \EE[\nu_{t+1}^{2}(x)]^\frac{1}{2}\bigg(  \frac{ \EE[\mu_{t+1}^{2}(x)]^\frac{1}{2}}{2^{M_t/2}}   + \sum_{s=t+1}^{T-1} \frac{L_{[t+1:s-1]}S_s\EE[(\mu_{s+1}^{2}(x))^2]^\frac{1}{2}}{2^{M_s+1}} \bigg).
    \end{align*}
\end{lemma}

\begin{proof}[Proof of Lemma~\ref{lemma:bias_bound_rtmlmc_grad}] 
We simplify notation by suppressing the dependence on~$x$ for all functions of~$x$ in this proof. In the following we will show that the conditional bias of the estimator~$\hat{G}_t(x)$ obeys the estimate
\begin{equation}
    \label{eq:bias-G_t}
    \begin{aligned}
    & \big\|\EE_{t-1}[\hat{G}_{t}]-G_{t} \big\|_2 \\ 
    & \qquad \leq \sum_{s=t}^{T-1}L_{[t:s-1]}S_s \EE_{t-1} [\nu_{s+1}^{2}]^\frac{1}{2}\bigg(\frac{ \EE_{t-1}[\mu_{s+1}^{2}]^\frac{1}{2}}{2^{M_s/2}} + \sum_{k=s+1}^{T-1} \frac{L_{[s+1:k-1]}S_k\EE_{t-1} [(\mu_{k+1}^{2})^2]^\frac{1}{2}}{2^{M_k+1}} \bigg)
    \end{aligned}
\end{equation}
for all~$t\in[T]$. As~$\EE[\hat{G}]=\EE[\hat{G}_1]$, the claim then follows immediately. We prove~\eqref{eq:bias-G_t} by backward induction on~$t$. As for the base case corresponding to~$t = T$, note that~$\hat{G}_T$ constitutes an unbiased estimator for~$G_T$. Thus, the left hand side of~\eqref{eq:bias-G_t} evaluates to~$0$, and the claim holds trivially (in fact, the right hand side of~\eqref{eq:bias-G_t} vanishes, too). As for the induction step corresponding to any~$t$ with~$t+1\in[T]$, assume that
\begin{align*}\label{eq:induction_hypothesis_bias_grad}
     \big\|\EE_{t}[\hat{G}_{t+1}]-G_{t+1} \big\|_2\leq \! \sum_{s=t+1}^{T-1}L_{[t+1:s-1]}S_s \EE_t[\nu_{s+1}^{2}]^\frac{1}{2}\bigg(\frac{ \EE_t[\mu_{s+1}^{2}]^\frac{1}{2}}{2^{M_s/2}}\! + \!\sum_{k=s+1}^{T-1} \!\frac{L_{[s+1:k-1]}S_k\EE_t[(\mu_{k+1}^{2})^2]^\frac{1}{2}}{2^{M_k+1}} \bigg).
\end{align*}
In addition, note that
\begin{equation}
    \begin{aligned}\label{eq:rtmlmc_grad_bias}
    \EE_{t-1}[\hat{G}_t] & =  \EE_{t-1}\left[\frac{1}{q_t(\lambda_t)} \left(\hat{g}_t^{\lambda_t}-\frac{1}{2}\hat{g}_t^{\lambda_t, {\rm e}} -\frac{1}{2}\hat{g}_t^{\lambda_t, {\rm o}} \right)\right] = \sum_{\ell=0}^{M_t} \EE_{t-1}\left[\hat{g}_t^{\ell}-\frac{1}{2}\hat{g}_t^{\ell, {\rm e}} -\frac{1}{2}\hat{g}_t^{\ell, {\rm o}} \right] \\
    &  = \EE_{t-1} [\hat{g}_t^{M_t}] 
    = \EE_{t-1}\Big[ \hat\EE_t^{M_t}[\hat{G}_{t+1}] \nabla_{x_t} f_t \big(\xi_t, \hat\EE_t^{M_t}[\hat{H}_{t+1}] \big)\Big],
\end{aligned}
\end{equation}
where the first equality follows from the definition of~$\hat{G}_t$ in~\eqref{eq:G_hat}, and the second uses the law of total expectation along with the independence of~$\lambda_t$ from all other random objects. The third equality holds because~$\hat{g}_t^{\ell,{\rm e}}$ and~$\hat{g}_t^{\ell,{\rm o}}$ share the distribution of~$\hat{g}_t^{\ell-1}$ conditional on $\mathcal F_t$ for all~$\ell \in [M_t]$ and because~$\hat{g}_t^{0,{\rm e}} = \hat{g}_t^{0,{\rm o}} = 0$. Finally, the fourth equality follows from Definition~\ref{def:rumlmc-grad}. Next, we introduce the shorthands
\begin{align*}
    A_1 & \coloneqq \Big\|\EE_{t-1}\Big[\hat\EE_t^{M_t}[\hat{G}_{t+1}] \left(\nabla_{x_t} f_t \big(\xi_t,\hat\EE_t^{M_t}[\hat{H}_{t+1}] \big) -\nabla_{x_t} f_t \left(\xi_t,F_{t+1}^T\right)\right)\Big] \Big\|_2\\
    A_2 & \coloneqq \Big\|\EE_{t-1}\Big[\Big(\EE_{t}[\hat{G}_{t+1}]-G_{t+1}\Big) \nabla_{x_t} f_t \left(\xi_t,F_{t+1}^T\right) \Big] \Big\|_2
\end{align*}
to simplify notation. Then, we have 
\begin{equation}\label{eq:bias_t_th_step}
\begin{aligned}
        \big\| \EE_{t-1}[\hat{G}_t] - G_t\big\|_2 & = \Big\|\EE_{t-1}\Big[\hat\EE_t^{M_t}[\hat{G}_{t+1}] \nabla_{x_t} f_t \big(\xi_t,\hat\EE_t^{M_t}[\hat{H}_{t+1}] \big) \Big] - \EE_{t-1}\Big[G_{t+1}\nabla_{x_t} f_t \big(\xi_t,F_{t+1}^T\big)\Big] \Big\|_2 \\
        &\leq \Big\|\EE_{t-1}\Big[\hat\EE_t^{M_t}[\hat{G}_{t+1}]\nabla_{x_t} f_t \big(\xi_t,\hat\EE_t^{M_t}[\hat{H}_{t+1}] \big) - \hat\EE_t^{M_t}[\hat{G}_{t+1}]\nabla_{x_t} f_t \big(\xi_t,F_{t+1}^T\big)\Big] \Big\|_2\\
        &\quad + \Big\|\EE_{t-1}\Big[\hat\EE_t^{M_t}[\hat{G}_{t+1}]\nabla_{x_t} f_t \big(\xi_t,F_{t+1}^T\big) - G_{t+1}  \nabla_{x_t} f_t \left(\xi_t,F_{t+1}^T\right) \Big]\Big\|_2 \! = A_1 \! + A_2,
\end{aligned}
\end{equation}
where the first equality exploits~\eqref{eq:rtmlmc_grad_bias} and the definition of~$G_t$ in~\eqref{eq:G_t}. The second equality follows from the law of total expectation and the elementary identity $\EE_{t}[\hat\EE_t^{M_t}[\cdot]]=\EE_t[\cdot]$. Next, we separately bound the terms~$A_1$ and~$A_2$. As for~$A_1$, we use the Cauchy-Schwarz and Jensen inequalities to obtain
\begin{align*}
        A_1 & \leq \EE_{t-1}\Big[\hat\EE_t^{M_t}\Big[\big\|\hat{G}_{t+1}\big\|_2\Big] \big\|\nabla_{x_t} f_t \big(\xi_t,\hat\EE_t^{M_t}[\hat{H}_{t+1}] \big) -\nabla_{x_t} f_t \left(\xi_t,F_{t+1}^T\right)\big\|_2 \Big] \\
        & \leq S_t \, \EE_{t-1}\Big[\big\|\hat{G}_{t+1}\big\|_2 \big\| \, \hat\EE_t^{M_t}[\hat{H}_{t+1}]  - F_{t+1}^T \big\|_2 \Big] \\
        & \leq S_t \, \EE_{t-1}\Big[\big\|\hat{G}_{t+1}\big\|_2^2\Big]^\frac{1}{2} \EE_{t-1}\Big[\big\|\hat\EE_t^{M_t}[\hat{H}_{t+1}]  - F_{t+1}^T \big\|_2^2 \Big]^\frac{1}{2} \\
        & = S_t \, \EE_{t-1}\Big[\big\|\hat{G}_{t+1}\big\|_2^2\Big]^\frac{1}{2}\Big( \EE_{t-1}\Big[\big\|\hat\EE_t^{M_t}[\hat{H}_{t+1}]  - \EE_t[\hat{H}_{t+1}] \big\|_2^2 \Big] 
        +\EE_{t-1}\Big[\big\|\EE_t[\hat{H}_{t+1}] - F_{t+1}^T\big\|_2^2 \Big]\Big)^\frac{1}{2}.
\end{align*}
Here, the second inequality follows from the smoothness of the integrand~$f_t$, as stipulated in Assumption~\ref{assump:smooth}. Note that we have also removed the first sample average $\hat\EE_t^{M_t}[\cdot]$. This is permissible because the second norm term does not depend on~$\xi_{t+1}$ and because the underlying samples are independent and identically distributed conditional on the information available at time~$t-1$. The third inequality and the equality follow from the Cauchy-Schwarz inequality and from a standard bias-variance decomposition of the mean squared error. 
By Lemma~\ref{lemma:bias_bound_rtmlmc_H_t}, the bias term admits the following upper bound.
\begin{align*}
    &\EE_{t-1}\Big[\big\|\EE_t[\hat{H}_{t+1}] - F_{t+1}^T\big\|_2^2 \Big] \leq \EE_{t-1}\bigg[\bigg(\sum_{s=t+1}^{T-1} \frac{L_{[t+1:s-1]} S_s \EE_{t}[\mu_{s+1}^2]}{2^{M_s + 1}} \bigg)^2\bigg]
\end{align*}
In addition, the variance term satisfies
\begin{align*}
    \EE_{t-1}\Big[\big\|\hat\EE_t^{M_t}[\hat{H}_{t+1}]  - \EE_t[\hat{H}_{t+1}] \big\|_2^2 \Big] = 2^{-M_t} \EE_{t-1}\Big[\big\|\hat{H}_{t+1}  - \EE_t[\hat{H}_{t+1}] \big\|_2^2 \Big] \leq 2^{-M_t} \EE_{t-1}[\mu_{t+1}^{2}],
\end{align*}
where the equality holds because the variance of the sample average $\hat\EE_t^{M_t}[\hat{H}_{t+1}]$ equals the variance of~$\hat{H}_{t+1}$ divided by the sample size~$2^{M_t}$. The inequality follows from the observation that the variance of~$\hat{H}_{t+1}$ is bounded below by its second moment and from the definition of~$\mu_{t+1}^{2}$ in~\eqref{eq:mu}. Substituting the bounds on these mean and the variance terms into the above bound on~$A_1$ then yields
\begin{align*}
    A_1 &\leq S_t \, \EE_{t-1}\Big[\big\|\hat{G}_{t+1}\big\|_2^2\Big]^\frac{1}{2} \Bigg( \EE_{t-1}\bigg[\bigg(\sum_{s=t+1}^{T-1} \frac{L_{[t+1:s-1]} S_s \EE_t[\mu_{s+1}^2]}{2^{M_s + 1}}\bigg)^2\bigg] + \frac{\EE_{t-1}[\mu_{t+1}^2]}{2^{M_t}} \Bigg)^\frac{1}{2}\\
    & \leq S_t \, \EE_{t-1}[\nu_{t+1}^{2}]^\frac{1}{2} \Bigg( \EE_{t-1}\bigg[\bigg(\sum_{s=t+1}^{T-1} \frac{L_{[t+1:s-1]} S_s \EE_t[\mu_{s+1}^2]}{2^{M_s + 1}}\bigg)^2\bigg]^\frac{1}{2} + \frac{\EE_{t-1}[\mu_{t+1}^{2}]^\frac{1}{2}}{2^{M_t/2}} \Bigg),
\end{align*}
where the second inequality holds because of the definition of~$\nu_{t+1}^2$ in~\eqref{eq:nu} and because $(a^2+b^2)^\frac{1}{2} \leq |a| + |b|$ for all~$a,b\in\mathbb R$. The Minkowski and Jensen inequalities further imply that
\begin{align*}
    \EE_{t-1}\bigg[\bigg(\sum_{s=t+1}^{T-1} \frac{L_{[t+1:s-1]} S_s \EE_t[\mu_{s+1}^{2}]}{2^{M_s + 1}}\bigg)^2\bigg]^\frac{1}{2} 
    &\leq \sum_{s=t+1}^{T-1}\EE_{t-1}\bigg[ \frac{L_{[t+1:s-1]}^2 S_s^2 \EE_t[\mu_{s+1}^{2}]^2 }{2^{2M_s + 2}}\bigg]^\frac{1}{2} \\
    &\leq \sum_{s=t+1}^{T-1}\frac{L_{[t+1:s-1]} S_s \EE_{t-1}[(\mu_{s+1}^{2})^2]^\frac{1}{2}}{2^{M_s + 1}}.
\end{align*}
Using this inequality to simplify that above bound on~$A_1$ yields
\begin{align*}
    A_1 \leq S_t \, \EE_{t-1}[\nu_{t+1}^{2}]^\frac{1}{2} \bigg( \sum_{s=t+1}^{T-1} \frac{L_{[t+1:s-1]} S_s \EE_{t-1}[(\mu_{s+1}^{2})^2]^\frac{1}{2}}{2^{M_s + 1}} + \frac{\EE_{t-1}[\mu_{t+1}^{2}]^\frac{1}{2}}{2^{M_t/2}} \bigg).
\end{align*}
As for the term~$A_2$, we use the Cauchy-Schwarz and Jensen inequalities to obtain
\begin{align*}
    &A_2 \leq  \EE_{t-1}\Big[\big\|\nabla_{x_t} f_t \left(\xi_t,F_{t+1}^T\right)\big\|_2 \, \big\|\EE_{t}[\hat{G}_{t+1}]-G_{t+1}  \big\|_2 \Big]\\
    &\quad \leq  L_t \; \EE_{t-1} \Big[\big\|\EE_{t}[\hat{G}_{t+1}]-G_{t+1} \big\|_2\Big]\\
    &\quad \leq L_t \;\EE_{t-1}\bigg[  \sum_{s=t+1}^{T-1}L_{[t+1:s-1]}S_s \EE_t[\nu_{s+1}^{2}]^\frac{1}{2} \bigg(\frac{\EE_{t}[\mu_{s+1}^{2}]^\frac{1}{2} }{2^{M_s/2}} + \bigg(\sum_{k=s+1}^{T-1} \frac{L_{[s+1:k-1]}S_k\EE_{t}[(\mu_{k+1}^{2})^2]^\frac{1}{2}}{2^{M_k + 1}}\bigg) \bigg) \bigg],
\end{align*}
where the second inequality follows from Assumption~\ref{assump:lipschitz}, and the third inequality exploits the induction hypothesis. Jensen’s inequality together with the law of total expectations further yields
\begin{align*}
&A_2 \leq  L_t \sum_{s=t+1}^{T-1}L_{[t+1:s-1]}S_s \EE_{t-1}[\nu_{s+1}^{2}]^\frac{1}{2} \bigg(\frac{\EE_{t-1}[\mu_{s+1}^{2}]^\frac{1}{2} }{2^{M_s/2}} + \bigg(\sum_{k=s+1}^{T-1} \frac{L_{[s+1:k-1]}S_k\EE_{t-1}[(\mu_{k+1}^{2})^2]^\frac{1}{2}}{2^{M_k + 1}}\bigg) \bigg) .
\end{align*}
Combining the bounds on~$A_1$ and~$A_2$ with~\eqref{eq:bias_t_th_step} and noting that $L_{[t:t-1]}=1$, we finally obtain
\begin{align*}
    &\big\| \EE_{t-1}[\hat{G}_t] - G_t\big\|_2\\
    &\quad \leq S_t \, \EE_{t-1}[\nu_{t+1}^{2}]^\frac{1}{2} \bigg(\frac{\EE_{t-1}[\mu_{t+1}^2]^\frac{1}{2}}{2^{M_t/2}} +\bigg(\sum_{s=t+1}^{T-1} \frac{L_{[t+1:s-1]} S_s \EE_{t-1}[(\mu_{s+1}^{2})^2]^\frac{1}{2}}{2^{M_s + 1}}\bigg) \bigg)   \\
    &\qquad\;\quad + L_t \sum_{s=t+1}^{T-1}L_{[t+1:s-1]}S_s \EE_{t-1}[\nu_{s+1}^{2}]^\frac{1}{2}\bigg(\frac{ \EE_{t-1}[\mu_{s+1}^{2}]^\frac{1}{2} }{2^{M_s/2}} + \bigg(\sum_{k=s+1}^{T-1} \frac{L_{[s+1:k-1]}S_k\EE_{t-1}[(\mu_{k+1}^{2})^2]^\frac{1}{2}}{2^{M_k + 1}}\bigg) \bigg)\\
    &\quad = \sum_{s=t}^{T-1} L_{[t:s-1]}S_s \EE_{t-1}[\nu_{s+1}^{2}]^\frac{1}{2} \bigg(\frac{ \EE_{t-1}[\mu_{s+1}^{2}]^\frac{1}{2}}{2^{M_s/2}} 
    + \sum_{k=s+1}^{T-1} \frac{L_{[s+1:k-1]}S_k\EE_{t-1}[(\mu_{k+1}^{2})^2]^\frac{1}{2}}{2^{M_k+1}}\bigg).
\end{align*}
This completes the induction step, and thus the claim follows.
\end{proof}

\end{document}
